\section{Implementation}

\subsection{Platform}

The sensing client is an iPhone. A lightweight activity-description model
runs locally on the phone and produces the draft caption without network
communication. The activity sentence bank and the embeddings required by the
draft stage are also stored locally.

Frames selected for verification are JPEG-encoded and sent through Wi-Fi to
a laptop edge server. The server is implemented using FastAPI and runs
SmolVLM-256M-Instruct as the verification model.

Before transmitting a full verification image, the iPhone sends a small
request containing the planned JPEG size and the confirmation deadline. The
laptop records its current CPU utilization and recent VLM execution latency,
predicts the verification time, computes the communication budget, and places
the request in the upload scheduler.

The scheduler returns an application-layer grant to the selected phone. The
phone then uploads the JPEG, the laptop performs VLM verification, and the
accept-or-rewrite result is returned to the phone.

This design requires no modification to the Wi-Fi MAC or access point and can
therefore be implemented entirely at the application layer.

\subsection{What is Already Running}

The current desktop prototype already implements the draft-and-verify logic.
For each frame it produces a draft caption, the top candidate sentences, the
confidence margin, the VLM output, and the accept-or-rewrite decision.

The next implementation step is to move the lightweight draft stage onto the
iPhone while retaining SmolVLM verification on the laptop. We will then add
timestamped Wi-Fi measurements and collect VLM latency traces under different
CPU loads and image configurations.

These traces will be used to train the edge-latency predictor and to evaluate
the compute-aware communication scheduler.

\subsection{Evaluation}

We will record short sessions of one person performing the behaviors included
in the activity sentence bank. The iPhone and laptop will operate on the same
Wi-Fi network.

For the AI pipeline, we will report:
\begin{itemize}
\item local draft latency;
\item fraction of drafts accepted and rewritten;
\item agreement between the local draft and VLM verification;
\item human assessment of the final caption on a small labeled subset.
\end{itemize}

For the system component, we will report:
\begin{itemize}
\item uplink, edge-compute, and downlink latency;
\item prediction error of $\hat T_i^{\mathrm{edge}}$;
\item end-to-end confirmation deadline-satisfaction ratio;
\item bytes uploaded per minute;
\item confirmation latency as the number of simultaneous clients increases.
\end{itemize}

We will evaluate three communication schedulers:
\begin{enumerate}
\item \emph{Round robin}: verification uploads are served without considering
their deadlines or compute latency;
\item \emph{Static budget}: communication deadlines are derived using a fixed
average VLM latency;
\item \emph{Predictive budget}: communication deadlines are dynamically
derived using the predicted VLM latency.
\end{enumerate}

Separately, the draft-and-verify pipeline will be compared with local draft
only and VLM verification on every frame. This allows the effect of the AI
pipeline and the effect of the communication scheduler to be evaluated
independently.

\subsection{Roles}

Both members will integrate the final demo, evaluation, and report.

One member will focus on the AI pipeline, including the local draft model,
activity sentence bank, acceptance rule, VLM verification, and caption-quality
evaluation.

The other member will focus on VLM latency profiling, edge-latency prediction,
communication-budget allocation, application-layer scheduling, and the
multi-client evaluation.

The draft-and-verify mechanism forms the information-processing core of the
project, while predictive deadline-aware communication forms the system and
networking core.